% =====================================================================
% VERSION ENSAE, REDUCTION SFE (phase 3, 1er octobre 2026) : chapitre 5, cascade dirigee entre piliers.
% Copie a texte constant de chapitres/07_cascade_dirigee.tex.
% Ne changent que les titres et les renvois repares, listes dans
% reduction/phase3.md.
% Lu par main_ensae.tex seul : ne jamais le faire lire par main_institut.tex,
% qui garde chapitres/.
% =====================================================================

\chapter{Cascade dirigée entre piliers}
\label{chap:cascade}

La matrice d'influence dirigée $W$ décrit la propagation entre les cinq piliers.
Le cadre latent ci-dessous sert à définir cette influence ; le capital est calculé par une marche
auto-évitante (\S\ref{sec:formalisation}). Le seuil latent $K_j$ est apparié à un seuil
statistique de matérialité $u_j$, non à une valeur fixée par DORA.

\section{Modèle de Vasicek à un facteur}

Le modèle structurel de Merton-Vasicek associe à chaque entité une variable latente $X$, qui
mesure le stress. Nous l'utilisons sous la probabilité historique : aucun prix ni aucune valeur de
firme n'est calculé. Le modèle s'écrit
\begin{equation}
X_i \;=\; \sqrt{\rho}\;Y \;+\; \sqrt{1-\rho}\;\varepsilon_i, \qquad \Var(X_i)=1,
\end{equation}
où $Y\sim\mathcal N(0,1)$ est le facteur systémique commun, $\varepsilon_i\sim\mathcal N(0,1)$ le
choc idiosyncratique, et $\rho\in[0,1]$ la part du systémique, commune à toutes les entités.
L'incident survient lorsque le stress franchit un seuil, $D_i=\mathbf 1\{X_i\ge K_i\}$, avec dans
le modèle classique, seuil de défaut du crédit, $K_i=\Phi^{-1}(1-\mathrm{PD}_i)$.

Deux hypothèses de ce modèle ne s'appliquent pas au cyber sous DORA. La dépendance y est
symétrique, avec un $Y$ et un $\rho$ uniques, alors que la matrice posée à dire d'expert est
nettement asymétrique : la norme de l'écart entre la matrice et sa transposée vaut $81\,\%$ de celle de la matrice. Le seuil $\Phi^{-1}(\mathrm{PD})$ suppose de plus une PD stable et mesurable et
une valeur de firme sous-jacente, qui n'existent pas en cyber.

% SOURCES-SCRIPTS: 01 03

\section{Facteur systémique par pilier}

Nous indexons désormais par entité~$i$ et pilier~$j\in\{1,\dots,5\}$, et remplaçons le facteur
unique par un facteur systémique par pilier $Y_j$ et une sensibilité propre $\rho_{ij}$ :
\begin{equation}
X_{ij} \;=\; \sqrt{\rho_{ij}}\;Y_j \;+\; \sqrt{1-\rho_{ij}}\;\varepsilon_{ij}.
\end{equation}
Les $Y_j$ sont corrélés entre eux par une matrice $\Sigma_Y$, qui décrit les dépendances croisées
au niveau systémique. Elle prolonge les modèles latents gaussiens du crédit (copule des temps de défaut, \citet{Li2000}).

Les données étudiées ne permettent pas de calibrer un $\rho_{ij}$ par entité. La probabilité marginale de non-conformité reste
à $0{,}35$, quelle que soit cette sensibilité ; seule la co-variation des cellules partageant un facteur la renseigne.
Même avec $Y_j$ connu, détecter un écart de $0{,}10$ demande environ $813$ périodes indépendantes,
d'après une approximation locale de la puissance à $80\,\%$.

Sur une coupe unique, deux fois la divergence de Kullback-Leibler atteint $0{,}47$ pour l'alternative la plus éloignée, $\rho_{ij}=0{,}99$.
Cette moyenne du double log-rapport de vraisemblance ne se compare pas au seuil de $3{,}84$.
Au niveau du pilier, il faut de l'ordre de $474$ périodes à quatre entités, et de $81$ sur une population entière.
Dans ce cadre, l'information vient de la répétition dans le temps.

Une extension serait le compromis hiérarchique $\rho_{ij}=\rho_j+u_{ij}$. Sa variante
$\mathrm{logit}\,\rho_{ij}=\mathrm{logit}\,\rho_j+u_{ij}$ reste dans $[0\,;\,1]$ par construction,
alors que le terme additif en sort avec une probabilité de $1{,}40\,\%$ pour une dispersion de
$0{,}20$. Le modèle publié garde une sensibilité unique et n'utilise pas cette extension.
Au chapitre~\ref{chap:multietats}, le facteur est celui de l'entité : l'approximation locale donne environ $374$ entités observées une fois chacune.

% SOURCES-SCRIPTS: 107

\section{Propagation dirigée}
\label{sec:baremes}

La propagation repose sur un modèle qualitatif de cascade, qui note les chaînes d'effondrement
entre piliers à partir de trois barèmes posés par l'auteur à dire d'expert
(figure~\ref{fig:baremes}). Aucune de leurs valeurs n'est mesurée. Elles sont ancrées sur la
structure du règlement, et nous n'en défendons que l'ordre.
\begin{itemize}[leftmargin=*,itemsep=2pt]
\item $\mathrm{ROOT}_p$ est la propension du pilier $p$ à amorcer une chaîne. Elle vaut $1{,}00$
  pour la gouvernance, qui fonde les autres contrôles, $0{,}90$ pour les tiers, surface d'attaque
  externe, et $0{,}30$ pour le partage d'informations, rarement déclencheur. Dans le moteur, un
  incident naît sur le pilier $p$ avec une probabilité proportionnelle à $\mathrm{ROOT}_p$, soit
  des parts d'amorce de $30\,\%$ (P1, gouvernance), $27\,\%$ (P4, tiers), $18\,\%$ (P2), $15\,\%$
  (P3) et $9\,\%$ (P5, partage d'informations).
\item $\mathrm{TRANS}_{jk}$ est la force avec laquelle la défaillance du pilier $j$ entraîne celle
  du pilier $k$. Ses vingt entrées vont de $0{,}10$ à $0{,}80$ et sont asymétriques : la
  gouvernance entraîne la gestion des incidents à $0{,}80$, l'inverse ne vaut que $0{,}20$. Ce
  sont des forces d'entraînement et non des probabilités, si bien que les lignes ne somment pas à
  un. Normalisée par l'équation~\eqref{eq:norm}, $\mathrm{TRANS}$ devient la matrice de contagion
  $W$.
\item $\mathrm{GBASE}_p$ est le score de gravité (au sens de l'AMDEC) du pilier $p$ s'il cède seul,
  sans rapport avec le gain de propagation $g$. Sur cette échelle ordinale, il vaut $7$ pour les tiers, qui transmettent la contagion, $6$ pour la gouvernance, $4$ et
  $3$ pour la gestion des incidents et les tests, $1$ pour le partage d'informations, qui est
  volontaire.
\end{itemize}
Dans le modèle qualitatif, le score de vraisemblance d'une chaîne ordonnée est le produit de $\mathrm{ROOT}$
au pilier d'amorce et de $\mathrm{TRANS}$ le long de ses maillons. Sa gravité part du
$\mathrm{GBASE}$ le plus élevé et croît avec les autres piliers touchés et avec la force moyenne des
maillons. Sa criticité, notée de $1$ à $10$, est la moyenne géométrique des deux. C'est elle que
mesure la dépendance à l'ordre du \S\ref{sec:verif-proprietes}. Le moteur de capital n'emploie que
$\mathrm{ROOT}$ et $\mathrm{TRANS}$. La sévérité y suit la même loi pour tous les piliers : c'est
une hypothèse, aucune base n'observant la sévérité par pilier, et les tests du
\S\ref{sec:profilage-variables} montrent seulement un effet faible des catégories de la base.

\begin{figure}[H]\centering
\figover{Z_baremes_expert.png}
\caption{Les trois barèmes d'expert du modèle qualitatif. (a)~$\mathrm{ROOT}$, propension de
chaque pilier à amorcer une chaîne. (b)~$\mathrm{GBASE}$, gravité du pilier s'il cède seul.
(c)~$\mathrm{TRANS}$, force avec laquelle le pilier source entraîne le pilier cible, asymétrique.}
\label{fig:baremes}
\end{figure}

Dans tout le rapport,
\[
  \boxed{\;W_{jk}\;=\;\text{propension du pilier } j \text{ (source) à entraîner la
  défaillance du pilier } k \text{ (cible).}\;}
\]
La ligne $j$ décrit ce que $j$ émet, la colonne $k$ ce que $k$ reçoit, comme dans le graphe
d'arêtes vivantes (définition~\ref{def:cascade}). Dans $\mathrm{TRANS}$, P1 a la plus grande somme
de ligne ($2{,}6$) et une somme de colonne trois fois inférieure à celle de P2 : il émet fort et
reçoit peu, et $W$ en garde les proportions. Le stress de $j$ subit celui des autres piliers de la même entité, et somme donc les colonnes :
\begin{equation}
X_{ij} \;=\; \underbrace{\sum_{k\neq j} W_{kj}\,X_{ik}}_{\text{contagion dirigée, } k \to j}
   \;+\; \sqrt{\rho_{ij}}\,Y_j \;+\; \sqrt{1-\rho_{ij}}\,\varepsilon_{ij},
\end{equation}
d'où, en forme matricielle sur les cinq piliers, la forme réduite
$\mathbf X_i=(I-W^{\top})^{-1}(\text{systémique}+\text{idio})$. Cette forme suppose
$\rho(W)<1$, le rayon spectral étant invariant par transposition. Si
$W=0$, on retombe sur le modèle à facteur par pilier, puis, avec un facteur unique, sur le Vasicek
standard. Les trois modèles latents s'emboîtent strictement. La figure~\ref{fig:reseau-W} représente $W$.

% SOURCES-SCRIPTS: 29 113

\section{Normalisation de Leontief et stabilité}

La matrice d'expert TRANS ne peut pas servir telle quelle : $(I-W)^{-1}$ ne représente
une cascade que si la série $I+W+W^2+\cdots$ converge, ce qui équivaut à $\rho(W)<1$. Pour TRANS,
on a $\rho(\mathrm{TRANS})=1{,}461>1$, et $(I-\mathrm{TRANS})^{-1}$ contient 21 entrées négatives
sur 25, si bien qu'un choc positif produirait un stress négatif ailleurs.

% FIGURE FIXEE ICI EN [H], ET ELLE EST APPELEE UNE PAGE PLUS HAUT (\S de la propagation
% dirigee). En [htbp] a son point d'appel elle derivait jusqu'a cette page-ci et s'inserait au
% milieu de la phrase de Leontief, coupant l'inverse B=(I-A)^{-1} de sa somme. Elle ne peut pas
% tenir sous son paragraphe d'appel (191 pt disponibles pour 242 demandes) et la retrecir
% mettrait ses etiquettes sous 5 pt. Posee ici, elle ouvre la page et le paragraphe la suit
% entier. Ne pas la remettre en [htbp].
\begin{figure}[H]\centering
\figover{H1_reseau_W.png}
\caption{Structure dirigée de $W$ sur les cinq piliers (P1 émet fort et reçoit peu) et
propagation d'un choc sur un sous-réseau.}
\label{fig:reseau-W}
\end{figure}

Le correctif vient du modèle entrées-sorties de Leontief, dont l'inverse
$B=(I-A)^{-1}=\sum_{k\ge0}A^k$ n'explose jamais parce que $A$ est une matrice de parts. Ses
colonnes somment à moins de un, la différence étant la valeur ajoutée, strictement positive.
Nous appliquons la même règle en lisant la ligne $j$ de $W$ comme les parts du stress que $j$
\emph{émet} vers les autres piliers :
\begin{equation}
\label{eq:norm}
\boxed{\;W \;=\; g \cdot \frac{\mathrm{TRANS}}{\max_j \sum_k \mathrm{TRANS}_{jk}}\;,
\qquad g \in (0,\,1]\;}
\end{equation}
où $g$, le gain de propagation, est la somme de ligne maximale de $W$ (P1, somme de ligne de TRANS $2{,}6$).
Le paramètre $g$ est posé, non estimé (valeurs $0{,}45$, $0{,}68$ et $0{,}90$ au
chapitre~\ref{chap:multietats}). $W$ est positive et chaque ligne somme à
au plus $g$, donc par Perron-Frobenius
\begin{equation}
\rho(W)\;\le\;\max_j\sum_k W_{jk}\;\le\;g\;\le\;1,
\end{equation}
avec $\rho(W)<1$ pour $g<1$. Numériquement, $\rho(W)=0{,}562\,g$, ce qui garde $\rho(W)<1$ à la borne
$g=1$ : $(I-W)^{-1}$ existe sur tout le domaine admissible. La
division par un même scalaire pour toutes les lignes préserve l'hétérogénéité de ce que reçoit
chaque pilier.

% SOURCES-SCRIPTS: 03 66

\section{Processus de branchement}
\label{sec:branch}

Les données ne permettent pas d'estimer le modèle simultané sur latentes. Nous le remplaçons par un
processus de branchement multitype : un incident sur le pilier~$j$ engendre, à la période suivante,
un nombre aléatoire d'incidents sur les autres piliers. La matrice de génération suivante $M$ compte des incidents :
\begin{equation}
M_{jk}\;=\;\Phi(\text{base}_k+W_{jk})-\Phi(\text{base}_k),
\end{equation}
où $\Phi(\text{base}_k)$ est l'incidence de fond du pilier cible.
Le rayon spectral $\rho(W)$ détermine le système latent simultané, tandis que $R_0=\rho(M)$ détermine
le branchement auxiliaire, dont $R_0<1$ garantit l'extinction avec une descendance moyenne finie. La cascade du modèle, qui ne revient jamais sur un pilier touché, s'arrête de toute façon. Avec une incidence de fond de
$4{,}5\,\%$ et $g=0{,}9$, valeurs posées, on obtient $\rho(W)=0{,}506$ et $R_0=0{,}054$. Les seuils critiques,
$\rho(W)=1$ à $g=1{,}78$ et $R_0=1$ à $g=8{,}1$, sont hors du domaine admissible $g\le1$
(figure~\ref{fig:branchement-R0-rapport}). Le seuil critique de la forme simultanée est plus de quatre fois plus proche que celui du branchement.

La progéniture ($(I-M)^{-1}$ moins l'identité) compte les incidents moyens causés par un incident sur chaque pilier.
Elle vaut $0{,}109$ (P1), $0{,}069$ (P4), $0{,}058$ (P2), $0{,}054$ (P3)
et $0{,}026$ (P5). Elle reproduit le classement ROOT (Spearman $=1{,}00$) à $g=0{,}9$.
ROOT et TRANS venant du même expert, cet accord est un test de cohérence interne.
Il ne détermine pas les valeurs de ROOT, qui restent posées séparément.

\begin{figure}[H]\centering
\figover{K3_branchement_R0.png}
\caption{Stabilité de la cascade et progéniture. (a)~Sur le domaine admissible $g\le1$, les deux
rayons spectraux restent inférieurs à $1$ ; les seuils critiques ($\rho(W)=1$ à $g=1{,}78$ ;
$R_0=1$ à $g=8{,}1$) sont hors domaine. (b)~La progéniture $(I-M)^{-1}$ reproduit le classement
ROOT du jugement d'expert.}
% SOURCES-SCRIPTS: 03
\label{fig:branchement-R0-rapport}
\end{figure}

% SOURCES-SCRIPTS: 03 29

\section{Seuil $K_j$, substitut de $\Phi^{-1}(\mathrm{PD})$}

Le seuil statistique $u_j$ se définit dans l'espace des sévérités observables. Le seuil $K_j$ s'applique à la variable latente, sans unité. Nous apparions leurs
probabilités, en choisissant $K_j$ pour que l'événement latent $\{X_{ij}\ge K_j\}$ ait la même
probabilité que l'incident majeur observable $\{S_{ij}\ge u_j\}$ :
\begin{equation}
\label{eq:Kj}
\boxed{\;K_j\;=\;F_{\mathrm{latent}}^{-1}(1-p_j), \qquad p_j=\mathbb P(S_{ij}\ge u_j).\;}
\end{equation}
CreditMetrics \citep{jpmorgan1997} inverse également une fréquence cumulée observée. Le règlement ne fixe pas $p_j$ : le seuil reste celui, statistique, de la méthode POT. $F_{\mathrm{latent}}$ décrit un stress sans unité, à ne pas confondre avec la sévérité GPD, exprimée en euros.

Le biais sur $K_j$ vient de la sous-déclaration. Si $q(s)$ est la
probabilité qu'un incident de sévérité $s$ soit déclaré, alors
$\hat p_j^{\text{naïf}}/p_j=\mathbb E[q(S)\mid S\ge u_j]<1$. Il en résulte une sous-estimation de
$p_j$, qui conduit à surestimer $K_j$ et à sous-provisionner. Ce biais s'effondre toutefois lorsque
le seuil de matérialité monte dans la queue (biais sur $K_j$, sans unité, de $+0{,}369$ à $u=0{,}5$~M\euro{} contre
$+0{,}001$ à $u=40$~M\euro{}). Nous modélisons la sévérité au-dessus de ce seuil par la théorie des valeurs
extrêmes. Dans une simulation où l'indice de queue vrai est connu ($\xi=0{,}9$), le biais médian
sur le capital est de $-0{,}8\,\%$ pour cette méthode, d'écart interquartile $0{,}62$ rapporté au capital vrai, contre
$-63{,}9\,\%$ pour la lognormale.

% SOURCES-SCRIPTS: 04 02

\section{Calibration et faisabilité}
\label{sec:calib-cascade}

Les sources examinées n'identifient pas l'asymétrie de $W$. Une chronologie fiable ou la loi complète des ensembles touchés pourraient renseigner la direction,
sous des hypothèses sur la cascade et les amorces. Nous testons un probit dynamique, adapté au bruit gaussien du modèle générateur.
Un logit gonflerait les magnitudes d'un facteur $2{,}1$.

Sur un panel de $800$ entités et $60$ périodes simulé par ce modèle, le probit retrouve $W$ dans ses unités (corrélation $0{,}905$, pente $1{,}085$).
La corrélation tombe à $0{,}38$ avec $50$ entités. Ce contrôle ne valide aucune calibration sur données réelles.
Le probit ne peut pas être estimé sur les sources disponibles (chapitre~\ref{ch:identifiabilite}), d'où le capital en bornes du chapitre~\ref{chap:partielle}.

% SOURCES-SCRIPTS: 03 05

\section{Définitions et propriétés formelles}
\label{sec:formalisation}

On note
$P=\{1,\dots,5\}$ les piliers, $W\in[0,1]^{P\times P}$ la matrice de contagion avec
$W_{jj}=0$, et $s_{\max}=\max_j\sum_k \mathrm{TRANS}_{jk}$. Les démonstrations sont
rassemblées en annexe~\ref{ann:cascade}.

\begin{definition}[Cascade indépendante]
\label{def:cascade}
À chaque couple ordonné $(j,k)$, $j\neq k$, on associe une variable de Bernoulli
indépendante $B_{jk}\sim\mathcal B(W_{jk})$. Le \emph{graphe d'arêtes vivantes} $G$ possède
l'arc $j\to k$ si et seulement si $B_{jk}=1$. Pour un pilier d'amorce $a$, l'\emph{ensemble
atteint} est
\[
  S_a \;=\; \{a\}\cup\{k\in P : \text{il existe un chemin dirigé } a\rightsquigarrow k \text{ dans } G\}.
\]
C'est la cascade indépendante standard \citep{kempe2003}, une percolation de liens sur le graphe pondéré par $W$.
\end{definition}

\begin{proposition}[Forme réduite et descendance espérée]
\label{prop:reduite}
La somme, sur les marches dirigées de longueur $n$ de $a$ vers $k$, du produit des poids de
leurs arêtes vaut $(W^n)_{ak}$, et la série de ces termes est
\[
  \sum_{n\ge 0}(W^n)_{ak} \;=\; \bigl((I-W)^{-1}\bigr)_{ak},
\]
convergente si et seulement si $\rho(W)<1$. Un pilier atteint l'étant par au moins un chemin vivant
sans pilier répété, on a la majoration $\mathbb E\,|S_a|\le \bigl[(I-W)^{-1}\mathbf 1\bigr]_a$, et
l'écart mesure les marches qui répètent un pilier et les piliers atteints par plusieurs chemins.
\end{proposition}

\begin{proposition}[Sous-criticité par normalisation de Leontief]
\label{prop:leontief}
Sous la normalisation~\eqref{eq:norm}, $W=(g/s_{\max})\,\mathrm{TRANS}$, la matrice $W$ est
positive et chacune de ses lignes somme à au plus $g$ ; par conséquent
\[
  \rho(W)\;\le\;\max_j\sum_k W_{jk}\;\le\;g\;\le\;1 .
\]
Pour $g<1$, cette borne garantit $\rho(W)<1$, donc la convergence de $(I-W)^{-1}$ et la sous-criticité de l'enveloppe de branchement. La cascade finie sans réactivation s'arrête indépendamment de cette condition. À la borne $g=1$ du domaine, elle ne donne que $\rho(W)\le 1$ : la stabilité tient
alors par la valeur calculée de $\rho(W)$, strictement inférieure à $1$.
\end{proposition}

\begin{proposition}[Loi exacte de l'ensemble atteint]
\label{prop:loi}
Pour tout $A\subseteq P$ avec $a\in A$,
\[
  \mathbb P(S_a=A)\;=\;R(A)\,B(A),\qquad
  B(A)=\!\!\prod_{j\in A,\;k\notin A}\!\!(1-W_{jk}),
\]
où $B(A)$ est la probabilité qu'aucune arête ne fuie hors de $A$, et $R(A)$, la probabilité
que $a$ atteigne \emph{tout} $A$ par les seules arêtes internes, vérifie $R(\{a\})=1$ et
\[
  R(A)\;=\;1-\sum_{\substack{A'\subsetneq A\\ a\in A'}} R(A')\!\!\prod_{j\in A',\;k\in A\setminus A'}\!\!(1-W_{jk}).
\]
\end{proposition}

\noindent Les capitaux des chapitres~\ref{chap:multietats} et~\ref{chap:resultats} sont calculés avec la variante à un successeur par pas, la marche
auto-évitante. Depuis le pilier $j$, la cascade passe au pilier $k$ avec la probabilité $W_{jk}$, et
s'arrête sinon, ou dès qu'elle revient sur un pilier touché. Les deux constructions ont la même
matrice moyenne et le même nombre moyen de descendants directs, et leurs lois se calculent
sans approximation. Les bornes du chapitre~\ref{chap:partielle} emploient la cascade indépendante, la
seule définie pour les matrices admissibles dont une ligne somme à plus de un. Passer de la marche au branchement déplace le besoin de capital du secteur de $-4\,\%$. Pour la marche, la
croissance du capital en $g$ est vérifiée numériquement au \S\ref{sec:invariance-g}.

\begin{proposition}[Monotonie stochastique]
\label{prop:mono}
Si $W\le W'$ terme à terme, alors $S_a(W)$ est stochastiquement dominé par $S_a(W')$. Toute
fonctionnelle croissante de $S_a$, en particulier la perte agrégée et sa VaR, est donc
croissante en chaque entrée de $W$. En revanche une perturbation \emph{antisymétrique}
$A$ déplace $W_{jk}$ et $W_{kj}$ en sens opposés : le capital n'est pas monotone en $A_{jk}$,
ce qui justifie de balayer les sommets de l'ensemble admissible
(chapitre~\ref{chap:partielle}) sans revendiquer de bornes exactes.
\end{proposition}

\begin{proposition}[Dépendance à l'ordre]
\label{prop:ordre}
Soit la vraisemblance d'une chaîne ordonnée
\[
\ell(i_1,\dots,i_n)=\mathrm{ROOT}_{i_1}\prod_{k=1}^{n-1}\mathrm{TRANS}_{i_k i_{k+1}}.
\]
Sur les chaînes de deux piliers, $\ell$ est invariante par
permutation si et seulement si
\[
\mathrm{ROOT}_i\,\mathrm{TRANS}_{ij}=\mathrm{ROOT}_j\,\mathrm{TRANS}_{ji}
\]
pour tout couple $(i,j)$. Cette condition est nécessaire à l'invariance pour toute longueur. Le barème la viole pour la paire (P1, P2), dont les deux ordres reçoivent des criticités différentes (chapitre~\ref{ch:positionnement}). Un score ou une copule des seules indicatrices ne représente pas l'ordre d'atteinte.
\end{proposition}

\begin{proposition}[Ce que la co-occurrence voit de la direction]
\label{prop:nonid}
Soit la décomposition unique $W=S+A$, avec $S=\tfrac12(W+W^\top)$ et $A=\tfrac12(W-W^\top)$.
\emph{(i)} Sous une amorce uniforme, de probabilité $r$ pour chaque pilier, la probabilité qu'un
même sinistre touche à la fois $j$ et $k$ vaut
\[
  \sum_{a\in P} r\,\mathbb P\bigl(\{j,k\}\subseteq S_a\bigr)
  \;=\; r\,(W_{jk}+W_{kj})+O\bigl(\lVert W\rVert^2\bigr)
  \;=\; 2r\,S_{jk}+O\bigl(\lVert W\rVert^2\bigr) :
\]
au premier ordre, la co-occurrence ne dépend de $W$ que par $S$ ; $W$ et $W^\top$ y sont
indiscernables. \emph{(ii)} Au-delà du premier ordre, ou sous une amorce non uniforme, la loi de
l'ensemble atteint dépend en général de $A$. $W$ et $W^\top$ ne sont donc pas observationnellement
équivalents pour la loi complète des ensembles de piliers touchés. \emph{(iii)} En général,
$\mathrm{SCR}(W)\neq\mathrm{SCR}(W^\top)$.
\end{proposition}

% SOURCES-SCRIPTS: 32 61 29 98

\section{Vérification numérique des propriétés}
\label{sec:verif-proprietes}

% SOURCES-SCRIPTS: 32
Sur les $26$ sous-ensembles de taille au moins deux, $22$ changent de criticité selon l'ordre,
l'écart atteignant $6$ points sur une échelle de $10$, c'est-à-dire $60\,\%$ de l'échelle. Sur la
paire $\{P_1,P_2\}$, $\mathrm{TRANS}[1][2]=0{,}80$ et $\mathrm{TRANS}[2][1]=0{,}20$ donnent des
criticités de $8$ et $2$ pour le même ensemble. Ni une copule ni un score additif ne peuvent
reproduire cet objet (proposition~\ref{prop:ordre}). La
non-transitivité de la dominance entre piliers, testée sous trois formulations, ne se vérifie pas,
avec zéro violation sur $420$ triplets, une valeur propre minimale de $+0{,}405$ et zéro cycle.
Nous ne la revendiquons donc pas.

Avec le classeur actuel, on a $\rho(W)=0{,}506$. La majoration $\E\,|S_a|\le[(I-W)^{-1}\mathbf 1]_a$
tient sur les cinq amorces et sur $300$ matrices de l'ensemble admissible. L'écart entre la borne
et la valeur exacte, de $10{,}6\,\%$ (amorce P5) à $19{,}9\,\%$ (amorce P1), comprend les
contributions des marches avec répétitions et des chemins multiples vers un même pilier. Nous gardons le calcul exact (proposition~\ref{prop:loi}).

Sur les transitions annuelles OpRisk, le test de permutation de l'asymétrie ne rejette pas $A=0$
(asymétrie $119$ contre $128 \pm 27$, $z=-0{,}33$, aucune arête significative sur $18$). Pour le
périmètre financier étudié, aux paramètres de base, retourner la matrice déplace pourtant le capital de $3{,}6\,\%$
($295$~M\euro) et la perte moyenne de $7{,}9\,\%$, et le pilier à remédier en premier passe de P1 à
P4. Sur l'ensemble admissible, l'écart entre une matrice et sa transposée atteint $1\,173$~M\euro{}.

La loi complète des ensembles touchés distingue les deux matrices, puisque, à $g=0{,}90$, cascade seule et hors choc commun P4, un sinistre touche en
moyenne $1{,}931$ pilier sous $W$ et $1{,}746$ sous $W^\top$. La part des sinistres multi-piliers
passe de $62{,}31\,\%$ à $50{,}77\,\%$, et la probabilité que P2 et P5 cèdent ensemble de $0{,}127$
à $0{,}066$. Sous amorce uniforme, l'écart n'est plus que de $0{,}86\,\%$ lorsque $W$ est ramenée
au cinquantième de sa valeur. Les statistiques disponibles ne séparent pas les deux matrices, dont les capitaux diffèrent pourtant.

% SOURCES-SCRIPTS: 40 03 98
